% Generated from lessons/0008-2-improper-integral-tests.html; edit the HTML source, then rerun the converter.
\section{反常积分的收敛判别法}
\lessonmark{反常积分的收敛判别法}

上一节会“算”反常积分；这一节要解决更常见的考试问题：算不出原函数时，怎样只判断它收敛、发散、绝对收敛还是条件收敛。

\subsection*{本节主线}

\begin{conceptbox}
\textbf{Cauchy 原理}
把“极限存在”翻译成尾部积分可以任意小。这是所有判别法的底层语言。
\end{conceptbox}

\begin{conceptbox}
\textbf{非负函数}
没有正负抵消，主要靠比较：与 \(p\)-积分、对数积分、幂函数比较。
\end{conceptbox}

\begin{conceptbox}
\textbf{变号函数}
先看绝对收敛；不绝对时，再试 Dirichlet 或 Abel，抓住振荡抵消。
\end{conceptbox}

\begin{notebox}
做题时别急着找原函数。反常积分题最常见的正确开局是：先找“坏点”，再找“比较对象”。
\end{notebox}

\subsection*{一、Cauchy 收敛原理}

以无穷区间为例，反常积分

\[
      \int_a^{+\infty}f(x)\,dx
    \]

收敛的充分必要条件是：对任意 \(\varepsilon>0\)，存在 \(A_0>a\)，使得只要 \(A',A''>A_0\)，就有

\[
      \left|\int_{A'}^{A''}f(x)\,dx\right|<\varepsilon.
    \]

这句话的直觉是：积分值要有极限，后面很远的尾巴就不能再贡献明显面积。

\subsubsection*{证明思路}

令 \(F(A)=\int_a^A f(x)\,dx\)。反常积分收敛就是 \(F(A)\) 当 \(A\to+\infty\) 有有限极限。把数列/函数极限的 Cauchy 条件套到 \(F(A)\) 上，就得到上式。

\subsection*{二、绝对收敛与条件收敛}

若

\[
      \int_a^{+\infty}|f(x)|\,dx
    \]

收敛，则称 \(\int_a^{+\infty}f(x)\,dx\) 绝对收敛。绝对收敛一定推出原积分收敛。

若 \(\int f\) 收敛，但 \(\int |f|\) 发散，则称条件收敛。

\begin{notebox}
考试顺序：先判 \(\int |f|\)。若绝对收敛，结束；若不绝对，再尝试 Dirichlet、Abel 或振荡替换。
\end{notebox}

\subsection*{三、非负函数的比较判别法}

\subsubsection*{定理 8.2.2：比较判别法}

设在 \([a,+\infty)\) 上有

\[
      0\le f(x)\le K\varphi(x).
    \]

则：

\begin{itemize}
 \item 若 \(\int_a^{+\infty}\varphi(x)\,dx\) 收敛，则 \(\int_a^{+\infty}f(x)\,dx\) 收敛；
 \item 若 \(\int_a^{+\infty}f(x)\,dx\) 发散，则 \(\int_a^{+\infty}\varphi(x)\,dx\) 发散。
 \end{itemize}

\subsubsection*{极限比较形式}

若 \(f(x),\varphi(x)\ge0\)，且

\[
      \lim_{x\to+\infty}\frac{f(x)}{\varphi(x)}=l,
    \]

则：

\begin{itemize}
 \item 若 \(0<l<+\infty\)，两个积分同敛散；
 \item 若 \(l=0\)，且 \(\int\varphi\) 收敛，则 \(\int f\) 收敛；
 \item 若 \(l=+\infty\)，且 \(\int\varphi\) 发散，则 \(\int f\) 发散。
 \end{itemize}

\subsubsection*{证明思路}

非负函数的积分上限函数 \(F(A)=\int_a^A f\) 单调递增。若能找到统一上界，就收敛；若被一个已发散的正函数从下方压住，就发散。

\subsection*{四、最常用的 \(p\)-型模板}

无穷远处：

\[
      \int_1^{+\infty}\frac{1}{x^p}\,dx
      \begin{cases}
      \text{收敛},&p>1,\\
      \text{发散},&p\le1.
      \end{cases}
    \]

有限端点瑕点：

\[
      \int_0^1\frac{1}{x^p}\,dx
      \begin{cases}
      \text{收敛},&p<1,\\
      \text{发散},&p\ge1.
      \end{cases}
    \]

对数模板：

\[
      \int_2^{+\infty}\frac{1}{x(\ln x)^p}\,dx
      \begin{cases}
      \text{收敛},&p>1,\\
      \text{发散},&p\le1.
      \end{cases}
    \]

\begin{notebox}
幂次决定主导行为；对数通常是“临界情况”才需要精细比较。
\end{notebox}

\subsection*{五、Abel 与 Dirichlet 判别法}

\subsubsection*{积分第二中值定理}

若 \(f\) 连续，\(g\) 单调，则存在 \(\xi\in[a,b]\)，使

\[
      \int_a^b f(x)g(x)\,dx
      =
      g(a)\int_a^\xi f(x)\,dx
      +g(b)\int_\xi^b f(x)\,dx.
    \]

它把“振荡函数 \(\times\) 单调函数”的积分拆成两个可控制的部分。

\subsubsection*{定理 8.2.5：Abel 判别法}

若 \(\int_a^{+\infty}f(x)\,dx\) 收敛，且 \(g(x)\) 在 \([a,+\infty)\) 上单调有界，则

\[
      \int_a^{+\infty}f(x)g(x)\,dx
    \]

收敛。

\subsubsection*{定理 8.2.5：Dirichlet 判别法}

若

\[
      F(A)=\int_a^A f(x)\,dx
    \]

在 \([a,+\infty)\) 上有界，且 \(g(x)\) 单调并满足 \(\lim_{x\to+\infty}g(x)=0\)，则

\[
      \int_a^{+\infty}f(x)g(x)\,dx
    \]

收敛。

\begin{center}
\small
\begin{tabularx}{\linewidth}{|>{\raggedright\arraybackslash}p{0.18\linewidth}|>{\raggedright\arraybackslash}p{0.42\linewidth}|>{\raggedright\arraybackslash}X|}
\hline
\textbf{看到什么} & \textbf{优先尝试} & \textbf{典型例子} \\ \hline
\(\sin x,\cos x\) 乘一个趋零因子 & Dirichlet & \(\int_1^\infty \frac{\sin x}{x}\,dx\) \\ \hline
已知 \(\int f\) 收敛，再乘单调有界函数 & Abel & \(\int_1^\infty \frac{\sin x\arctan x}{x}\,dx\) \\ \hline
非负、无振荡 & 比较或极限比较 & \(\int_1^\infty \frac{\ln x}{x^p}\,dx\) \\ \hline
\end{tabularx}
\end{center}

\subsection*{六、无界函数反常积分的对应版本}

若 \(f\) 在 \(x=b\) 附近无界，Cauchy 原理变成：对任意 \(\varepsilon>0\)，存在 \(\delta>0\)，使任意 \(\eta',\eta''\in(0,\delta)\) 都有

\[
      \left|\int_{b-\eta'}^{b-\eta''}f(x)\,dx\right|<\varepsilon.
    \]

比较判别法也完全平行。例如若 \(x\to b-\) 时

\[
      f(x)\le \frac{K}{(b-x)^p},
    \]

则 \(p<1\) 时可推出收敛；若 \(f(x)\ge \frac{K}{(b-x)^p}\)，则 \(p\ge1\) 时可推出发散。

\begin{notebox}
无穷远和有限瑕点只是“坏点位置”不同：一个看 \(x\to+\infty\)，一个看 \(x\to b-\) 或 \(x\to a+\)。判别逻辑相同。
\end{notebox}

\subsection*{七、考试判别流程}

\begin{enumerate}
 \item 先找坏点：无穷远、端点无界、内点无界。
 \item 若有多个坏点，必须拆开，每段分别判别。
 \item 若被积函数非负，优先比较 \(x^{-p}\)、\((b-x)^{-p}\)、\([x(\ln x)^p]^{-1}\)。
 \item 若有 \(\sin,\cos\) 振荡，先判绝对收敛；绝对不行再试 Dirichlet 或 Abel。
 \item 若参数题，分别检查每个坏点给出的参数限制，最后取交集。
 \end{enumerate}

\subsection*{八、即时判断练习}

\begin{quickcheck}
\(\int_1^\infty \frac{\sin x}{x}\,dx\) 的性质是？

\tcblower
\textbf{答案：}b
\end{quickcheck}

\begin{quickcheck}
若 \(0\le f(x)\sim \frac{1}{x^{3/2}}\)，则 \(\int_1^\infty f(x)\,dx\)？

\tcblower
\textbf{答案：}a
\end{quickcheck}

\begin{quickcheck}
\(\int_0^1 \frac{1}{x^p}\,dx\) 收敛的条件是？

\tcblower
\textbf{答案：}c
\end{quickcheck}

\subsection*{九、课后题训练}

下面按教材第 324-325 页习题 1-16 整理。题面尽量完整保留，提示只给解题入口；写作业时建议先遮住提示独立判断一遍。

\begin{exercisebox}
\textbf{习题 1 \badge{比较判别法}}

(1) 证明比较判别法（定理 8.2.2）；

(2) 举例说明，当比较判别法的极限形式中 \(l=0\) 或 \(+\infty\) 时，\(\int_a^\infty\varphi(x)\,dx\) 和 \(\int_a^\infty f(x)\,dx\) 的敛散性可以产生各种不同的情况。

\begin{hintbox}
(1) 用非负积分上限函数单调性；(2) 选 \(x^{-p}\) 做样本，调整 \(p\) 跨过 1。
\end{hintbox}
\end{exercisebox}

\begin{exercisebox}
\textbf{习题 2 \badge{Cauchy 判别法}}

证明 Cauchy 判别法及其极限形式（定理 8.2.3 及其推论）。

\begin{hintbox}
把 \(K/x^p\) 作为比较函数；极限形式通过“充分远处夹在常数倍之间”转化为普通比较。
\end{hintbox}
\end{exercisebox}

\begin{exercisebox}
\textbf{习题 3 \badge{非负函数敛散性}}

讨论下列非负函数反常积分的敛散性：

\[
      \begin{aligned}
      &(1)\ \int_1^{+\infty}\frac{1}{\sqrt{x^3-e^{-2x}+\ln x+1}}\,dx,\\
      &(2)\ \int_1^{+\infty}\frac{\arctan x}{1+x^3}\,dx,\\
      &(3)\ \int_0^{+\infty}\frac{1}{1+x|\sin x|}\,dx,\\
      &(4)\ \int_1^{+\infty}\frac{x^q}{1+x^p}\,dx\quad(p,q\in\mathbb R^+).
      \end{aligned}
      \]

\begin{hintbox}
(1)(2) 看无穷远主项；(3) 注意不能被 \(\frac1x\) 直接控制，尝试在 \(|\sin x|\) 小的区间上估计；(4) 无穷远等价于 \(x^{q-p}\)。
\end{hintbox}
\end{exercisebox}

\begin{exercisebox}
\textbf{习题 4 \badge{Cauchy 主值}}

证明：对非负函数 \(f(x)\)，\((cpv)\int_{-\infty}^{+\infty}f(x)\,dx\) 收敛与 \(\int_{-\infty}^{+\infty}f(x)\,dx\) 收敛是等价的。

\begin{hintbox}
非负时没有正负抵消，主值的对称截断若有界，就能推出两侧普通积分都有界。
\end{hintbox}
\end{exercisebox}

\begin{exercisebox}
\textbf{习题 5 \badge{绝对/条件/发散}}

讨论下列反常积分的敛散性（包括绝对收敛、条件收敛和发散）：

\[
      \begin{aligned}
      &(1)\ \int_2^{+\infty}\frac{\ln\ln x}{\ln x}\sin x\,dx,\\
      &(2)\ \int_1^{+\infty}\frac{\sin x}{x^p}\,dx\quad(p\in\mathbb R^+),\\
      &(3)\ \int_1^{+\infty}\frac{\sin x\arctan x}{x^p}\,dx\quad(p\in\mathbb R^+),\\
      &(4)\ \int_0^{+\infty}\sin(x^2)\,dx,\\
      &(5)\ \int_a^{+\infty}\frac{p_m(x)}{q_n(x)}\sin x\,dx.
      \end{aligned}
      \]

其中 \(p_m(x)\) 和 \(q_n(x)\) 分别是 \(m\) 和 \(n\) 次多项式，且 \(q_n(x)\) 在 \(x\in[a,+\infty)\) 范围内无零点。

\begin{hintbox}
核心是“振荡 \(\sin x\)”配合单调趋零因子。第 (5) 先比较有理函数的无穷阶 \(x^{m-n}\)。
\end{hintbox}
\end{exercisebox}

\begin{exercisebox}
\textbf{习题 6 \badge{瑕积分判别法}}

设 \(f(x)\) 在 \([a,b]\) 只有一个奇点 \(x=b\)，证明定理 8.2.3' 和定理 8.2.5'。

\begin{hintbox}
把无穷远版本中的 \(x\to+\infty\) 换成 \(x\to b-\)，所有“尾部”改成“靠近瑕点的小段”。
\end{hintbox}
\end{exercisebox}

\begin{exercisebox}
\textbf{习题 7 \badge{有限区间瑕积分}}

讨论下列非负函数反常积分的敛散性：

\[
      \begin{aligned}
      &(1)\ \int_0^1\frac{1}{\sqrt[3]{x^2(1-x)}}\,dx,\\
      &(2)\ \int_0^1\frac{\ln x}{x^2-1}\,dx,\\
      &(3)\ \int_0^{\pi/2}\frac{1}{\cos^2x\,\sin^2x}\,dx,\\
      &(4)\ \int_0^{\pi/2}\frac{1-\cos x}{x^2}\,dx,\\
      &(5)\ \int_0^1|\ln x|^p\,dx,\\
      &(6)\ \int_0^1x^{p-1}(1-x)^{q-1}\,dx,\\
      &(7)\ \int_0^1x^{p-1}(1-x)^{q-1}|\ln x|\,dx.
      \end{aligned}
      \]

\begin{hintbox}
逐个找端点坏点。第 (6)(7) 是 Beta 型，分别在 0 和 1 附近给出参数限制。
\end{hintbox}
\end{exercisebox}

\begin{exercisebox}
\textbf{习题 8 \badge{参数型敛散性}}

讨论下列反常积分的敛散性：

\[
      \begin{aligned}
      &(1)\ \int_0^{+\infty}\frac{x^{p-1}-x^{q-1}}{\ln x}\,dx\quad(p,q\in\mathbb R^+),\\
      &(2)\ \int_0^{+\infty}\frac{1}{\sqrt[3]{x(x-1)^2(x-2)}}\,dx,\\
      &(3)\ \int_0^{+\infty}\frac{\ln(1+x)}{x^p}\,dx,\\
      &(4)\ \int_0^{+\infty}\frac{\arctan x}{x^p}\,dx,\\
      &(5)\ \int_0^{\pi/2}\frac{\sqrt{\tan x}}{x^p}\,dx,\\
      &(6)\ \int_0^{+\infty}x^{p-1}e^{-x}\,dx,\\
      &(7)\ \int_0^{+\infty}\frac{1}{x^p+x^q}\,dx,\\
      &(8)\ \int_2^{+\infty}\frac{1}{x^p\ln^q x}\,dx.
      \end{aligned}
      \]

\begin{hintbox}
这是本节最值得反复练的一组。对每题列出所有坏点，再分别比较。参数条件最后取交集。
\end{hintbox}
\end{exercisebox}

\begin{exercisebox}
\textbf{习题 9 \badge{综合振荡题}}

讨论下列反常积分的敛散性：

\[
      \begin{aligned}
      &(1)\ \int_0^{+\infty}\frac{x^{p-1}}{1+x^q}\,dx,\\
      &(2)\ \int_1^{+\infty}\frac{x^p\sin x}{1+x^q}\,dx\quad(p\ge0),\\
      &(3)\ \int_0^{+\infty}\frac{e^{\sin x}\cos x}{x^p}\,dx,\\
      &(4)\ \int_0^{+\infty}\frac{e^{\sin x}\sin 2x}{x^p}\,dx,\\
      &(5)\ \int_0^1\frac{1}{x^p}\cos\frac{1}{x^2}\,dx,\\
      &(6)\ \int_1^{+\infty}\frac{\sin\left(x+\frac1x\right)}{x^p}\,dx\quad(p>0).
      \end{aligned}
      \]

\begin{hintbox}
第 (3)(4) 可利用 \(d(e^{\sin x})=e^{\sin x}\cos x\,dx\) 以及周期抵消；第 (5) 常用 \(t=1/x\) 或 \(t=1/x^2\)。
\end{hintbox}
\end{exercisebox}

\begin{exercisebox}
\textbf{习题 10 \badge{振荡复合}}

证明反常积分

\[
        \int_0^{+\infty}x\sin(x^4)\sin x\,dx
      \]

收敛。

\begin{hintbox}
尝试分部积分或利用快速振荡 \(\sin(x^4)\)；关键是不要只用绝对值估计。
\end{hintbox}
\end{exercisebox}

\begin{exercisebox}
\textbf{习题 11 \badge{绝对收敛推出平方收敛}}

设 \(\int_a^{+\infty}f(x)\,dx\) 绝对收敛，且 \(\lim_{x\to+\infty}f(x)=0\)。证明 \(\int_a^{+\infty}f^2(x)\,dx\) 收敛。

\begin{hintbox}
充分远处 \(|f(x)|\le1\)，所以 \(f^2(x)\le |f(x)|\)。
\end{hintbox}
\end{exercisebox}

\begin{exercisebox}
\textbf{习题 12 \badge{单调无界的必要条件}}

设 \(f(x)\) 单调，且当 \(x\to0+\) 时 \(f(x)\to+\infty\)。证明：\(\int_0^1f(x)\,dx\) 收敛的必要条件是

\[
        \lim_{x\to0+}x f(x)=0.
      \]

\begin{hintbox}
用单调性在区间 \([x,2x]\) 上给出面积下界。
\end{hintbox}
\end{exercisebox}

\begin{exercisebox}
\textbf{习题 13 \badge{单调乘子极限}}

设 \(\int_a^{+\infty}f(x)\,dx\) 收敛，且 \(xf(x)\) 在 \([a,+\infty)\) 上单调减少。证明：

\[
        \lim_{x\to+\infty}x(\ln x)f(x)=0.
      \]

\begin{hintbox}
类比习题 12，用 \([x,x^2]\) 或相邻倍增区间估计尾部积分。
\end{hintbox}
\end{exercisebox}

\begin{exercisebox}
\textbf{习题 14 \badge{导函数振荡积分}}

设 \(f(x)\) 单调下降，且 \(\lim_{x\to+\infty}f(x)=0\)。证明：若 \(f'(x)\) 在 \([0,+\infty)\) 上连续，则反常积分

\[
        \int_0^{+\infty}f'(x)\sin^2x\,dx
      \]

收敛。

\begin{hintbox}
写 \(\sin^2x=\frac12-\frac12\cos2x\)。一部分直接由 \(f(\infty)-f(0)\) 控制，另一部分用 Dirichlet/分部积分。
\end{hintbox}
\end{exercisebox}

\begin{exercisebox}
\textbf{习题 15 \badge{平方可积}}

若 \(\int_a^{+\infty}f^2(x)\,dx\) 收敛，则称 \(f(x)\) 在 \([a,+\infty)\) 上平方可积。讨论：

(1) 对两种反常积分分别探讨 \(f(x)\) 平方可积与 \(f(x)\) 的反常积分收敛之间的关系；

(2) 对无穷区间的反常积分，举例说明平方可积与绝对收敛互不包含；

(3) 对无界函数的反常积分，证明平方可积必定绝对收敛，但逆命题不成立。

\begin{hintbox}
无穷区间和有限区间的结论不同。有限区间可用 Cauchy-Schwarz；无穷区间要找反例。
\end{hintbox}
\end{exercisebox}

\begin{exercisebox}
\textbf{习题 16 \badge{参数综合压轴}}

证明反常积分

\[
        \int_1^{+\infty}\frac{\sin x}{x^p+\sin x}\,dx
      \]

当 \(p\le\frac12\) 时发散，当 \(\frac12<p\le1\) 时条件收敛，当 \(p>1\) 时绝对收敛。

\begin{hintbox}
把分母按 \(x^p\) 提出来：\(\frac{\sin x}{x^p+\sin x}=\frac{\sin x}{x^p}\cdot\frac{1}{1+\sin x/x^p}\)，再展开或比较主项。决定条件收敛和发散的关键是二阶项 \(\sin^2x/x^{2p}\)。
\end{hintbox}
\end{exercisebox}

主参考：陈纪修、于崇华、金路《数学分析》第三版上册，第八章 §2，教材第 315-325 页。

上一节：第八章 §1：反常积分的概念和计算。如果某个判别法或某道题卡住，直接在浏览器里选中对应段落问我。
